\section{模块五：频域图像增强，从一行像素开始}

\noindent\textbf{这一模块要解决什么？}\quad 一幅图里可能同时有大范围的明暗变化、物体边界、草地纹理和重复条纹。我们希望把其中某些变化保留下来，把另一些变化减弱。\textbf{频域（frequency domain）}把问题改写成：“变化有多快？沿哪个方向变化？这一类变化有多强？”先理解这三个问题，再学公式，滤波器就不再只是需要背诵的字母。

本模块沿用原讲义模块五及原题 B4、D6、D7、E6 的编号。条纹小例子、四点变换、边界补零示例等是为理解原题补充的必要背景；它们不是原讲义中额外承诺的考试题目。原稿中不严谨的结论会在相应位置明确订正。

\subsection{先看一行亮度：什么叫变化快，什么叫变化慢}

把灰度图像的一行像素依次抄下来，就得到一串数。数值大表示亮，数值小表示暗。例如下面每行都有八个像素：
\[
\begin{array}{c|rrrrrrrr}
\text{位置 }x&0&1&2&3&4&5&6&7\\\hline
\text{几乎不变}&100&100&100&100&100&100&100&100\\
\text{缓慢变亮}&40&50&60&70&80&90&100&110\\
\text{反复明暗}&40&110&40&110&40&110&40&110
\end{array}
\]
第一行没有随位置发生的明暗变化；第二行要走过较长距离才有明显变化；第三行只走一个像素就从暗变亮或从亮变暗。频域中的“高”和“低”，首先描述这种\textbf{空间变化的快慢}，不描述像素灰度数值的高低。一整块灰度为 240 的白墙仍可以主要是低频；在 20 和 30 之间密集交替的暗条纹也可以含很高的空间频率。

\textbf{空间频率（spatial frequency）}是每单位距离内完成的周期数。一个完整周期可以是“亮到暗，再回到亮”。若相邻两条亮纹之间相隔 $T$ 个像素，则
\[
\nu=\frac{1}{T}.
\]
这里 $T$ 是\textbf{周期（period）}，单位为像素；$\nu$ 是空间频率，单位为周期/像素。这个公式用于把可见的条纹间距转换成变化快慢：$T=8$ 时 $\nu=1/8=0.125$ 周期/像素；$T=4$ 时 $\nu=1/4=0.25$ 周期/像素。后一组条纹更密，所以频率更高。若距离单位换成毫米，频率单位就相应成为周期/毫米。

\begin{center}
\begin{tikzpicture}[x=0.68cm,y=0.46cm,>=Latex,font=\small]
\draw[->] (0,0)--(16.8,0) node[right]{$x$};
\draw[->] (0,-.2)--(0,3.5) node[above]{灰度};
\draw[dashed,gray] (0,1.6)--(16,1.6);
\draw[thick,domain=0:16,samples=150] plot(\x,{1.6+1.1*cos(45*\x)});
\draw[<->] (0,3)--(8,3) node[midway,fill=white]{一个周期：8 像素};
\node[anchor=west] (a) at (0,-1){慢变化：16 像素中出现 2 个周期};
\begin{scope}[yshift=-4.4cm]
\draw[->] (0,0)--(16.8,0) node[right]{$x$};
\draw[->] (0,-.2)--(0,3.5) node[above]{灰度};
\draw[dashed,gray] (0,1.6)--(16,1.6);
\draw[thick,domain=0:16,samples=180] plot(\x,{1.6+1.1*cos(90*\x)});
\draw[<->] (0,3)--(4,3) node[midway,fill=white]{4 像素};
\node[anchor=west] at (0,-1){快变化：16 像素中出现 4 个周期};
\end{scope}
\end{tikzpicture}
\end{center}

\textbf{读图检查：}两条曲线的平均高度和上下摆动幅度相同，改变的是每段距离中摆动的次数。这说明“频率高”既不等于“平均更亮”，也不等于“对比度更大”。有限图像的边缘和裁切也会影响频谱；“缓慢渐变主要对应低频”是对主体变化的描述，不是保证其他频率完全为零。

\subsection{用波拼出亮度：幅度、相位和直流}

\subsubsection{一条最简单的波要用哪些量描述}
先补一个几何背景：在半径为 1 的圆上，从最右侧出发逆时针转过角度 $\theta$，到达点的横坐标叫 $\cos\theta$（余弦），纵坐标叫 $\sin\theta$（正弦）。转一整圈回到起点，所以这两个数会周期性重复，适合描述明暗反复出现的变化。\textbf{弧度（radian）}是角度的另一种单位：一整圈 $360^\circ=2\pi$ 弧度，四分之一圈 $90^\circ=\pi/2$ 弧度。公式里的角度默认按弧度读。

\begin{center}
\begin{tabular}{@{}lllll@{}}
\toprule
转过的角度&弧度&到达位置&$\cos\theta$&$\sin\theta$\\
\midrule
$0^\circ$&$0$&最右侧&$1$&$0$\\
$90^\circ$&$\pi/2$&最上方&$0$&$1$\\
$180^\circ$&$\pi$&最左侧&$-1$&$0$\\
$270^\circ$&$3\pi/2$&最下方&$0$&$-1$\\
$360^\circ$&$2\pi$&回到最右侧&$1$&$0$\\
\bottomrule
\end{tabular}
\end{center}

只看横坐标，就按 $1,0,-1,0,1$ 变化。把这个上下起伏的数乘上一个幅度，再加上基准亮度，就得到可见的明暗波。一个余弦亮度变化可以写为
\[
f(x)=b+A\cos\!\left(2\pi\nu x+\phi\right).
\]
这里 $f(x)$ 是位置 $x$ 的灰度，$b$ 是基准亮度，$A\geq0$ 是\textbf{振幅（amplitude）}，控制上下摆动的大小；$\nu$ 是频率；$\phi$ 是\textbf{相位（phase）}，控制波峰从哪里开始；$\pi$ 是圆周率。$2\pi$ 表示一个完整周期的角度，因而 $\nu x$ 先给出走过了多少周期，再乘 $2\pi$ 转成余弦函数使用的角度。

若把这条波直接用作不截断的 8 位非负灰度，需要 $b-A\geq0$ 且 $b+A\leq255$，保证最暗和最亮处都在可显示范围内。下面例子满足这一条件。

例如取 $b=100$、$A=20$、$\nu=1/4$、$\phi=0$，得到
\[
f(x)=100+20\cos(\pi x/2).
\]
逐点代入 $x=0,1,2,3$：余弦值分别为 $1,0,-1,0$，灰度就是 $120,100,80,100$。到 $x=4$ 又变成 120，因此每四个像素重复一次。

如果只把 $A$ 改为 40，结果为 $140,100,60,100$，明暗差更大，周期仍是四个像素。如果只把 $\phi$ 改为 $\pi/2$，结果为 $100,80,100,120$，只是波峰移动了位置。\textbf{频率决定重复速度，幅度决定这一成分有多强，相位决定它如何对齐。}

\subsubsection{复杂图像为什么可以谈“有哪些频率”}
一条亮度曲线通常不是单独一条波，但可以由多条不同频率、幅度和相位的波相加得到。\textbf{傅里叶分析（Fourier analysis）}就是找出这些波的“配方”。在有限离散数据上，离散傅里叶变换提供一套可逆的配方：不丢掉系数时，能够用逆变换重建原数据。

例如
\[
f(x)=100+20\cos(2\pi x/8)+5\cos(2\pi x/2).
\]
第一项 100 是处处一样的基准；第二项每八个像素一个周期，振幅 20；第三项每两个像素一个周期，振幅 5。以 $x=0$ 为例，$f(0)=100+20+5=125$；以 $x=2$ 为例，$f(2)=100+20\cos(\pi/2)+5\cos(2\pi)=105$。这里高频成分更弱，但确实存在。滤去第三项，就去掉了细密起伏，而保留较缓慢的明暗变化。

“把图像分解成波”不是说真实场景由许多条纹实体组成，而是换一种数学方式描述同一组像素。类似一个位置可以用经纬度表示，也可以换一套坐标表示；改变表示方式不等于改变图像。

\subsubsection{直流为什么与平均亮度有关}
\textbf{直流分量（direct-current component，DC）}指零频率成分。零频率的波不再随位置起伏，所以对应一个常数。采用后文的正变换不归一化、逆变换除以像素数的约定时，
\[
F(0,0)=\sum_{x=0}^{M-1}\sum_{y=0}^{N-1} f(x,y),
\qquad
\bar f=\frac{F(0,0)}{MN}.
\]
$M,N$ 分别是图像在两个坐标方向的像素数，$MN$ 是总像素数，$\bar f$ 是平均灰度。因此\textbf{直流系数等于总和，除以总像素数才是平均值}。例如 $2\times2$ 图像的四个像素全为 10，总和为 40，故 $F(0,0)=40$，平均值为 $40/4=10$。原稿把直流直接说成“平均灰度”，是在说它代表的含义；数值上须看变换的归一化约定。

\subsection{认识二维频谱：一个点到底代表什么}

\subsubsection{先约定坐标，避免把横竖记反}
本模块统一把 $x$ 看作竖直方向的行坐标，把 $y$ 看作水平方向的列坐标；$u$ 与 $x$ 配对，$v$ 与 $y$ 配对。图像尺寸为 $M$ 行、$N$ 列。一组二维条纹可写成
\[
f(x,y)=b+A\cos\!\left[2\pi\left(\frac{k_x x}{M}+\frac{k_y y}{N}\right)+\phi\right].
\]
$k_x$ 表示沿整幅图的竖直跨度走过多少周期，$k_y$ 表示沿水平跨度走过多少周期。$k_x/M$ 和 $k_y/N$ 是两个方向的空间频率。若 $k_x=0,k_y=4$，则只沿水平方向变化；沿竖直方向保持不变，所以看到的是\textbf{竖条纹}。对应频率点出现在频谱中心的水平方向上。

\begin{center}
\begin{tikzpicture}[font=\small,>=Latex]
\begin{scope}
\fill[gray!15] (0,0) rectangle(4,3);
\foreach \a in {0,...,79}{\pgfmathsetmacro{\shadeval}{50+40*cos(18*\a)}\fill[black!\shadeval] ({0.05*\a},0) rectangle({0.05*\a+0.051},3);}
\draw (0,0) rectangle(4,3);
\draw[->] (0,-.35)--(4,-.35) node[right]{$y$};
\draw[->] (-.35,3)--(-.35,0) node[below]{$x$};
\node[align=center] at (2,-1){竖条纹：向右移动时明暗交替\\向下移动时不变};
\end{scope}
\draw[->,thick] (4.6,1.5)--(5.8,1.5) node[midway,above]{变换};
\begin{scope}[xshift=7.8cm,yshift=1.5cm]
\draw (-2,-1.5) rectangle(2,1.5);
\draw[->] (-1.8,0)--(2.35,0) node[right]{$\Delta v$};
\draw[->] (0,-1.3)--(0,1.95) node[above]{$\Delta u$};
\fill (0,0) circle(2.5pt);
\fill (-1.2,0) circle(3pt);\fill (1.2,0) circle(3pt);
\node[above] at (-1.2,0){$-k_y$};\node[above] at (1.2,0){$+k_y$};
\node[below right] at (0,0){DC};
\node[align=center] at (0,-2.5){频率点沿水平方向分开\\两点连线垂直于条纹走向};
\end{scope}
\end{tikzpicture}
\end{center}

本图为方便读坐标把频谱竖轴正向画向上；软件把矩阵行号画向下时，竖轴方向可能相反，但“哪个方向发生变化”“连线垂直于条纹”不变。横条纹沿竖直方向发生变化，亮点沿频谱竖直方向分开。斜条纹也遵循同一个规则：沿条纹走，亮度变化少；垂直穿过条纹，亮暗交替最多。

\subsubsection{频谱位置、亮度和相位不能混为一谈}
对每个频率位置 $(u,v)$，变换得到一个通常为复数的系数 $F(u,v)$。复数可以写成
\[
F(u,v)=a(u,v)+j\,b(u,v),\qquad j^2=-1.
\]
$a$ 和 $b$ 是这个系数的实部、虚部；这里的 $b(u,v)$ 与前面例子中的基准常数 $b$ 不同。二者是记录波的强弱与对齐位置的一种方便写法。我们最常画出的量是
\[
|F(u,v)|=\sqrt{a(u,v)^2+b(u,v)^2},
\qquad
\theta(u,v)=\operatorname{atan2}(b(u,v),a(u,v)).
\]
$|F|$ 是\textbf{幅度谱（magnitude spectrum）}；$\theta$ 是\textbf{相位谱（phase spectrum）}；$\operatorname{atan2}$ 根据两部分的符号确定正确角度。例如 $F=3+4j$ 时幅度为 $\sqrt{9+16}=5$，相位约为 $53.13^\circ$。若 $F$ 改为 $-3-4j$，幅度仍为 5，但相位不同，重建出的波峰位置也不同。

因此，频谱图中“某处很亮”的意思是\textbf{该频率的成分很强}，不是原图同一位置有个很亮的像素。空间位置与频谱位置是两套坐标；一个频率系数通常影响整幅图像，不能把右上角的频谱点直接解释成原图右上角的物体。

幅度的数值范围常很大，直接显示时强直流会掩盖弱分量。常用显示式为
\[
s(u,v)=c_{\mathrm{disp}}\ln\bigl(1+|F(u,v)|\bigr).
\]
$s$ 是显示用的亮度，$c_{\mathrm{disp}}>0$ 是把结果放进屏幕灰度范围的缩放系数；加 1 避免对零取对数。比如幅度 0 和 99 分别变为 $\ln 1=0$ 和 $\ln100\approx4.605$，大幅度被压缩。\textbf{这只是为了看清频谱，真正滤波仍使用原来的复系数 $F$，不能拿显示图 $s$ 代替它。}

\subsubsection{为什么实图像会出现成对频率点}
一幅普通灰度图的像素是实数。它的频谱具有\textbf{共轭对称性（conjugate symmetry）}：
\[
F((-u)\bmod M,(-v)\bmod N)=\overline{F(u,v)}.
\]
$\overline{F}$ 表示把复数虚部的符号反过来；“$\bmod$”表示把索引绕回合法范围。例如 $3+4j$ 的共轭为 $3-4j$，两者幅度相同。所以中心化后的幅度谱通常关于中心成对对称。

不必先掌握复数技巧也能理解原因：一条实余弦波可由正、负两个相反频率的复波配对表示。两项的虚部互相抵消，留下实数。对非零且非特殊边界频率，单一理想余弦会产生一对亮点。DC 和某些奈奎斯特频率点可与自己配对，因此“所有亮点都另有一个不同点”并非普遍成立。

图像平移主要改变各个波的相位。在\textbf{循环平移、没有新裁切}的离散模型下，幅度谱不变；现实中移动后又裁掉一块内容，幅度也可能改变。只保存幅度而丢弃相位，一般不能唯一恢复原图。

图像整体旋转时，频谱也随之旋转相同角度：原来水平的变化方向旋转后，频率点连线也转向新方向。这是连续理想模型的性质；离散图像做插值旋转或重新裁切时，幅度可能还受到插值与边界的影响。

\subsection{DFT 与 IDFT：把“配方”写成可计算的公式}

\subsubsection{正变换是在问：每一类波用了多少}
还需要认识复指数这个简写。\textbf{欧拉公式（Euler's formula）}为
\[
e^{j\theta}=\cos\theta+j\sin\theta,\qquad
e^{-j\theta}=\cos\theta-j\sin\theta.
\]
$j$ 是满足 $j^2=-1$ 的虚数单位；$\theta$ 是弧度角。这个写法把一对余弦、正弦合在一个数里，方便同时记录波的强弱和对齐位置，并不意味着像素本身必须是复数。利用前面的四分之一圈表，$\theta=0,\pi/2,\pi,3\pi/2$ 时，$e^{-j\theta}$ 依次等于 $1,-j,-1,j$。例如 $\theta=\pi/2$ 时，$\cos\theta=0,\sin\theta=1$，所以 $e^{-j\pi/2}=0-j\times1=-j$。

\textbf{离散傅里叶变换（discrete Fourier transform，DFT）}把像素数组变成频率系数数组。本模块采用
\[
F(u,v)=\sum_{x=0}^{M-1}\sum_{y=0}^{N-1}
 f(x,y)\exp\!\left[-j2\pi\left(\frac{ux}{M}+\frac{vy}{N}\right)\right].
\]
逐一读符号：$f(x,y)$ 是原图的像素；$x=0,\ldots,M-1$ 是行号，$y=0,\ldots,N-1$ 是列号；$u=0,\ldots,M-1$ 和 $v=0,\ldots,N-1$ 是两方向的频率索引。$\sum$ 表示把所有像素的贡献加起来；$\exp(t)=e^t$ 是指数函数；$j$ 是虚数单位。指数中的 $ux/M+vy/N$ 负责指定一组波，负号约定它是正变换。

这条公式并不要求在做滤波题时逐像素手算。它告诉我们三件事：一个频率系数由整幅图决定；系数含幅度与相位；当 $u=v=0$ 时指数为 0，$e^0=1$，因而该系数就是像素总和。

\textbf{快速傅里叶变换（fast Fourier transform，FFT）}是快速计算 DFT 的算法，不是一种不同的变换。类似用更快的乘法算法算出同一个乘积，FFT 与直接求和在数学上得到同样的 DFT，浮点舍入误差除外。

\subsubsection{逆变换是在按配方把所有波加回去}
\textbf{逆离散傅里叶变换（inverse DFT，IDFT）}为
\[
f(x,y)=\frac{1}{MN}\sum_{u=0}^{M-1}\sum_{v=0}^{N-1}
F(u,v)\exp\!\left[j2\pi\left(\frac{ux}{M}+\frac{vy}{N}\right)\right].
\]
变量含义不变，但现在遍历所有频率，把各个成分加到位置 $(x,y)$。与正变换相比，指数符号变成正号，并除以总像素数 $MN$。正、逆两式必须成套使用；别的教材可能把归一化系数放在正变换或拆到两边，只要前后一致，结论相同。

用前面的 $2\times2$ 常数图验证：只有 $F(0,0)=40$ 非零，任意位置的逆变换为 $(1/4)\times40\times1=10$，正好重建原灰度。这就是“直流代表整体亮度”的一次实际运算。

\subsubsection{四个像素的完整小例子}
为看清正、逆变换在做什么，先退回一维，取四个像素
\[
f=[12,10,8,10].
\]
其一维 DFT 是 $F(k)=\sum_{x=0}^{3}f(x)e^{-j2\pi kx/4}$，其中 $k=0,1,2,3$ 是四个频率索引。$k=0$ 时四个乘数都是 1，所以 $F(0)=12+10+8+10=40$。

当 $k=1$ 时，依次使用乘数 $1,-j,-1,j$，于是
\[
F(1)=12-10j-8+10j=4.
\]
$k=2$ 时乘数是 $1,-1,1,-1$，故 $F(2)=12-10+8-10=0$。$k=3$ 时与 $k=1$ 共轭，得到 $F(3)=4$。所以
\[
F=[40,4,0,4].
\]
这个配方对应 $f(x)=10+2\cos(2\pi x/4)$：常数为 10，一条振幅为 2 的余弦；它在频谱中表现为成对的两个系数 4。\textbf{DFT 系数的幅度并不直接等于原余弦的振幅}，还受样本数和正负频率配对影响；本例每个系数为 $4\times2/2=4$。

逆变换在 $x=0$ 处为 $(40+4+0+4)/4=12$；在 $x=2$ 处为 $(40-4+0-4)/4=8$。这样就能看到“频率配方”怎样重新变成像素，而不是把 DFT 只当作一个神秘缩写。

\subsection{滤波的核心：给每种频率乘一个权重}

\textbf{频率响应或传递函数（frequency response/transfer function）}$H(u,v)$ 给每一个频率位置指定一个权重。频域滤波的核心仅是
\[
G(u,v)=H(u,v)F(u,v).
\]
$F$ 是输入频谱，$H$ 是滤波器，$G$ 是输出频谱。若某位置 $H=1$，该分量不变；$H=0$，该分量完全去掉；$H=0.2$，只剩原幅度的五分之一；$H=2$，幅度加倍。对本模块常用的实数、非负 $H$，相位保持不变；一般复数 $H$ 也可以改变相位。

例如某频率系数为 $F=3+4j$，原幅度为 5。若 $H=0.4$，则 $G=1.2+1.6j$，幅度为 $\sqrt{1.2^2+1.6^2}=2$，正好是原来的 0.4 倍。

在四像素例子中，取 $H=[1,0.5,0,0.5]$，得到
\[
G=[40,2,0,2],\qquad g=[11,10,9,10].
\]
平均亮度仍为 10，原来上下摆动 2，现在上下摆动 1。若改成 $H=[1,0,0,0]$，则输出处处为 10，只剩平均亮度。若反过来只去掉直流，$H=[0,1,1,1]$，输出为 $[2,0,-2,0]$。出现负值是正常的：这是亮度\textbf{变化量}，不能直接当成已经放进显示范围的图像。

\subsubsection{与空间卷积的关系及适用条件}
\textbf{卷积定理（convolution theorem）}说，在匹配的变换和边界约定下，空间卷积对应频谱逐点相乘：
\[
g=f*h\quad\Longleftrightarrow\quad G=FH.
\]
这里 $*$ 表示卷积，$h$ 是空间域滤波核，$H$ 是 $h$ 的傅里叶变换。以一维为例，$g(x)=\sum_t h(t)f(x-t)$；$t$ 枚举核的位置，每一项是权重乘邻近像素，再求和。

这个对应关系适用于固定核的\textbf{线性、平移不变}滤波。中值滤波需要先排序再选中间数，通常不能写成一个固定 $H$ 乘所有输入的频谱；同态滤波在取对数后的域里可用线性滤波，但整个原灰度输入到输出过程一般不是线性的。不能把“任何图像增强”都说成与某个固定频域乘法严格等价。

\subsection{为什么频谱要中心化：搬动索引，不是创造细节}

未经移动的 DFT 把 DC 放在数组左上角 $(0,0)$。低频不仅出现在左上角附近，也会绕到其他角落，肉眼较难看出“中心到外围”的规律。\textbf{频谱中心化（spectrum centering）}把零频率搬到图中央，便于显示和画圆形滤波器。

有两条常用、必须保持前后一致的路线：
\begin{enumerate}[label=\arabic*.]
\item 先算 DFT，再对频谱执行 \texttt{fftshift} 重排，把零频率搬到中央；在中心化频谱上乘 $H$ 后，用 \texttt{ifftshift} 搬回原顺序，再做 IDFT。
\item 当变换网格的两个尺寸都是偶数时，在正变换前把图像乘 $(-1)^{x+y}$；在中心化频谱上乘 $H$，做 IDFT 后再乘一次 $(-1)^{x+y}$，还原空间图像。
\end{enumerate}

$(-1)^{x+y}$ 是一张正负交替的棋盘：当 $x+y$ 为偶数时为 $+1$，为奇数时为 $-1$。
\[
\begin{bmatrix}
1&-1&1&-1\\-1&1&-1&1\\1&-1&1&-1\\-1&1&-1&1
\end{bmatrix}.
\]
它不是把图像整体变暗，而是利用傅里叶变换的调制性质，使频率索引沿两个方向各平移半个网格。对于偶数 $P\times Q$ 的变换网格，DC 移到 $(P/2,Q/2)$。再次乘棋盘能还原，因为 $[(-1)^{x+y}]^2=1$。

\textbf{为什么强调偶数尺寸？}当尺寸为奇数时，“半个网格”不是整数索引，这个棋盘调制不能简单等同于软件的整数格 \texttt{fftshift}。一般实现可采用成套的 \texttt{fftshift}/\texttt{ifftshift}，并按实际中心索引建立 $H$。中心化是方便，不是必需条件：也可以始终用未中心化坐标设计 $H$，只要滤波器与频谱的索引位置一致。

\textbf{最常见的错误：}对未中心化的 $F$，直接乘一个画在数组中央的低通圆盘。这样保留的会是错误位置的频率。另一个错误是只移动一次，逆变换前没有移回，结果带上交替符号。记住：\textbf{显示顺序可以改，配对关系不能错。}

\subsection{补零与边界：为什么右边会“绕回”左边}

\subsubsection{DFT 看待边界的方式}
有限长 DFT 把数组理解为一个周期的内容，周期结束后重新从开头接上。二维情况可以想象把同一张图片向四周反复平铺。因此在原尺寸网格上直接进行 DFT 相乘再 IDFT，得到的是\textbf{循环卷积（circular convolution）}：右边可能绕回左边，下边也可能绕回上边。

真实照片的左边和右边往往不连续，周期接缝会产生人为突变。滤波时这些接缝可能污染边界。因此“边界处理”既有数学上的卷积长度问题，也有图像外部应该取什么值的建模问题。

\subsubsection{用四则运算看懂循环卷积与线性卷积}
令一维图像 $f=[1,2,3]$，滤波核 $h=[1,1]$。把数组外面看成零，\textbf{线性卷积（linear convolution）}逐项计算为
\[
\begin{aligned}
g(0)&=1\times1=1,\\
g(1)&=2\times1+1\times1=3,\\
g(2)&=3\times1+2\times1=5,\\
g(3)&=3\times1=3.
\end{aligned}
\]
得到 $g=[1,3,5,3]$，长度为 $3+2-1=4$。如果只用长度 3 的 DFT，最后的结果没有位置存放，$g(3)=3$ 会绕回并加到 $g(0)$，得到循环卷积 $[4,3,5]$。这不是数值误差，而是计算了另一种边界模型。

正确的补零做法是把两者都放在至少长度 4 的共同网格：$f_p=[1,2,3,0]$，$h_p=[1,1,0,0]$。用四点 DFT 相乘再逆变换，就可以完整容纳 $[1,3,5,3]$，不发生首尾折叠。下标 $p$ 表示经过填充的数组。

\subsubsection{二维尺寸应怎样确定}
若图像有 $M\times N$ 个像素，有限滤波核有 $K\times L$ 个系数，那么完整线性卷积的尺寸为
\[
(M+K-1)\times(N+L-1).
\]
为了用 DFT 计算这个线性卷积，选共同变换尺寸 $P\times Q$ 满足
\[
P\geq M+K-1,\qquad Q\geq N+L-1,
\]
并把\textbf{图像和核都}填到这个共同尺寸。$P,Q$ 可以更大，以方便快速算法；不是必须取二的整数次幂。例如 $256\times256$ 图像与 $3\times3$ 核，至少需要 $258\times258$ 的共同网格；只把核补到 $256\times256$ 仍不足以容纳完整线性卷积。

此外，核的原点必须安排正确。以核左上角为索引零直接做完整卷积，随后应按核中心偏移选取与输入对齐的部分；若把核中心作为原点，则要按软件约定将正、负位移正确放入网格。\textbf{补零解决重叠，原点与裁剪解决对齐，两者不是一件事。}

\subsubsection{直接画频域滤波器，也不能跳过边界问题}
原稿称“若 $H$ 直接在频域构造，如 ILPF，则无此问题”，这不成立。只要在有限 DFT 网格上做乘法，逆变换仍对应这个网格上的循环卷积。理想低通的空间响应延伸很远，更不能因未显式写出 $h$ 就忽略边界。

若没有一个指定的有限核，便不能机械套用某个有限 $K,L$ 获得严格的全局等价。常见做法是扩展图像、在更大的网格上滤波，再裁取关心的区域；必要时结合镜像延拓等边界模型。补零可减少绕回污染，但零与边界像素之间也可能形成突变。应说明采用的是零延拓、镜像、复制边界还是周期延拓。

还有一种不同用途：为了\textbf{让频谱显示的采样更密}而给图像补零。它使频谱曲线看起来更平滑，却没有增加原图的细节信息，也没有凭空提高对两个邻近真实频率的分辨能力。必须分清“卷积不折叠”和“频谱插值更密”这两种目的。

\subsection{用半径选频率：三种低通滤波器}

\subsubsection{先把频谱位置换成到中心的距离}
对已中心化的偶数 $P\times Q$ 变换网格，定义
\[
\Delta u=u-P/2,\qquad \Delta v=v-Q/2,
\qquad D(u,v)=\sqrt{(\Delta u)^2+(\Delta v)^2}.
\]
$\Delta u,\Delta v$ 是相对中心的有符号偏移；$D$ 是到中心的欧氏距离；$D_0>0$ 是滤波器的尺度或约定截止半径。原题 E6 没有额外补零，因此网格记成 $M=N=256$；本模块用 $P,Q$ 区分原图尺寸和变换尺寸。

例如偏移 $(3,4)$ 的点，$D=\sqrt{3^2+4^2}=5$。以中心画半径 5 的圆，圆上所有方向都受到相同处理，这叫\textbf{径向对称}。上述 $D$ 的单位是“频率格”，并非直接的周期/毫米。如果图像两个方向尺寸或像素物理间距不同，要用归一化的物理频率计算圆形距离，不能把不同单位直接相加。

\textbf{低通滤波（low-pass filtering）}让慢变化通过，压低快变化。结果通常更平滑，但也可能模糊边缘、纹理和小目标。它不知道“哪种细节重要”，只按频率做选择。

\subsubsection{理想低通 ILPF：一刀切}
\textbf{理想低通滤波器（ideal low-pass filter，ILPF）}为
\[
H_{\mathrm I}(u,v)=
\begin{cases}
1,&D(u,v)\leq D_0,\\
0,&D(u,v)>D_0.
\end{cases}
\]
$H_{\mathrm I}$ 是理想低通的权重；$D_0$ 内全部保留，外面全部丢弃。本书约定边界 $D=D_0$ 取 1，采用别的边界约定时须说明。例如 $D_0=10$，$D=8$ 时 $H=1$，$D=12$ 时 $H=0$。优点是选频范围清楚；代价是过于突然的频率截断会引起后面解释的振铃。

原稿的“物理不可实现”需要限定：连续、无限范围的精确理想频率响应，不能由有限长度的空间核精确实现；但\textbf{有限 DFT 网格上的零一圆盘当然可以计算}。不能把它理解为计算机无法使用理想低通。

\subsubsection{巴特沃斯低通 BLPF：用阶数控制坡度}
\textbf{巴特沃斯低通滤波器（Butterworth low-pass filter，BLPF）}在本讲义中定义为
\[
H_{\mathrm B}(u,v)=\frac{1}{1+\left[D(u,v)/D_0\right]^{2n}},\qquad n\geq1.
\]
$n$ 是正整数阶数，控制过渡带有多陡；$D/D_0$ 是无单位的相对距离；指数是 $2n$。例如 $D_0=10$、$n=2$、$D=5$ 时，
\[
H_{\mathrm B}=\frac{1}{1+(5/10)^4}
=\frac{1}{1+1/16}=\frac{16}{17}\approx0.9412.
\]
所以这个分量保留约 $94.12\%$ 的振幅。若 $D=20$，则 $H=1/(1+2^4)=1/17\approx0.05882$，仅保留约 $5.88\%$。当 $D=D_0$ 时比值为 1，任何阶数都得到 $H=1/(1+1)=0.5$。

$n$ 越大，$D_0$ 内更接近 1，外部更接近 0，形状越接近理想低通。原稿建议的 $n=2$ 至 $3$ 可以作为尝试的起点，而非通用最优值。低阶通常较平缓，但不能无条件承诺“所有低阶实现绝无振铃”；滤波器参数、采样和边界处理都会影响实际结果。

\subsubsection{高斯低通 GLPF：平缓下降}
\textbf{高斯低通滤波器（Gaussian low-pass filter，GLPF）}为
\[
H_{\mathrm G}(u,v)=\exp\!\left[-\frac{D(u,v)^2}{2D_0^2}\right].
\]
$D_0$ 在这个写法中是频域的尺度参数。$D=0$ 时指数为 0，$H=1$；$D=D_0$ 时 $H=e^{-1/2}\approx0.60653$，不是 0.5；$D=2D_0$ 时 $H=e^{-2}\approx0.13534$。因此三种滤波器使用同一个 $D_0$，并不表示它们在该位置有相同的衰减。

\begin{center}
\begin{tikzpicture}[x=3.5cm,y=3.2cm,>=Latex,font=\small]
\draw[->] (0,0)--(2.75,0) node[right]{$D/D_0$};
\draw[->] (0,0)--(0,1.22) node[above]{$H$};
\draw[dotted] (1,0)--(1,1.08);
\draw[dotted] (0,.5)--(2.6,.5);
\draw[very thick] (0,1)--(1,1)--(1,0)--(2.6,0);
\draw[thick,dashed,domain=0:2.6,samples=100] plot(\x,{1/(1+\x^4)});
\draw[thick,dash dot,domain=0:2.6,samples=100] plot(\x,{exp(-\x^2/2)});
\foreach \x in {1,2}{\draw (\x,0)--(\x,-.025) node[below]{$\x$};}
\node[left] at (0,1){1};\node[left] at (0,.5){0.5};
\node[anchor=west,fill=white] at (1.4,1.08){实线：理想};
\node[anchor=west,fill=white] at (1.4,.94){虚线：巴特沃斯 $n=2$};
\node[anchor=west,fill=white] at (1.4,.80){点划线：高斯};
\end{tikzpicture}
\end{center}

这张图只画了从频谱中心往外走的一条半径，不是原图的灰度剖面。横轴往右表示变化更快；纵轴越低表示这种变化保留得越少。$D_0$ 变小，更多频率被压低，通常更模糊；$D_0$ 变大，更多细节得以保留，去噪也可能变弱。

\subsubsection{怎样选择半径：能量比例只是一个参考}
幅度的平方 $|F(u,v)|^2$ 常用来度量该频率的能量贡献。若希望了解半径 $D_0$ 内包含多少频谱能量，可计算
\[
\eta(D_0)=\frac{\displaystyle\sum_{D(u,v)\leq D_0}|F(u,v)|^2}
{\displaystyle\sum_{u,v}|F(u,v)|^2}.
\]
$\eta$ 为圆内能量占总能量的比例，分子只累加圆内频率，分母累加全部频率；分母为零的全零图像需单独处理。例如总能量为 1000，圆内为 920，则 $\eta=920/1000=0.92=92\%$。

原稿提出以包含约 $90\%$ 至 $95\%$ 能量的圆作选择参考，这可作为尝试参数的一种办法，\textbf{不是普遍最优标准}。若 DC 极强，它可能单独占绝大多数能量，极小半径也满足比例，却把重要边缘全部抹掉。实际还要结合目的、噪声分布与视觉结果；有时需单独分析去掉 DC 后的能量分布，并清楚说明采用的定义。

\subsubsection{纠错：0.5 到底是负 3 dB 还是负 6 dB}
\textbf{分贝（decibel，dB）}表达比值时，必须先说清是振幅比还是功率比。对振幅传递函数 $H$，
\[
A_{\mathrm{dB}}=20\log_{10}|H|.
\]
$A_{\mathrm{dB}}$ 是振幅增益的分贝数，$\log_{10}$ 是以 10 为底的对数。将 $H=0.5$ 代入，
\[
20\log_{10}(0.5)\approx-6.0206\ \mathrm{dB}.
\]
对应功率比为 $|H|^2=0.25$，用功率公式 $10\log_{10}(0.25)$ 仍得到 $-6.0206$ dB。\textbf{功率比为 0.5} 才对应 $10\log_{10}(0.5)\approx-3.0103$ dB，其振幅比为 $\sqrt{0.5}\approx0.7071$。

有些工程教材定义巴特沃斯\textbf{振幅}为 $1/\sqrt{1+(D/D_0)^{2n}}$，那时 $D_0$ 处确实约为 $-3$ dB。但原题使用的是\textbf{没有平方根}的 $1/[1+(D/D_0)^{2n}]$。不能保留这一公式，却搬用另一种定义的截止结论。

\subsection{为什么边界会出现明暗晕圈：振铃不是“还有噪声”}

\textbf{振铃（ringing）}是边缘附近出现交替的过亮、过暗波纹。原图即使完全没有噪声，经过陡峭截止的滤波器，也可能出现振铃。因此不能看见波纹就断言去噪没有去干净。

频域与空间域通过傅里叶变换联系：频域截断越突然，对应的空间核越容易长距离振荡。对一维理想低通，空间核是 sinc 型；二维\textbf{圆形}理想低通的空间核严格说是带贝塞尔函数的 jinc 型径向函数。两者都具有旁瓣且包含负值。把二维圆盘直接称为“二维 sinc”只是一种不精确的类比。

不必先计算 sinc，也可以用一个有负系数的小核看清后果。考虑从 0 突然跳到 1 的阶跃：$f(x)=0$（$x<0$），$f(x)=1$（$x\geq0$）。取仅用于示意的对称核
\[
h(-2)=-0.1,\quad h(-1)=0.2,\quad h(0)=0.8,
\quad h(1)=0.2,\quad h(2)=-0.1.
\]
系数总和为 1，能保持远离边缘处的常数；但负系数可能把输出拉到原范围之外。按 $g(x)=\sum_{t=-2}^{2}h(t)f(x-t)$，得到
\[
\begin{array}{c|rrrrrrr}
x&-3&-2&-1&0&1&2&3\\\hline
f(x)&0&0&0&1&1&1&1\\
g(x)&0&-0.1&0.1&0.9&1.1&1&1
\end{array}
\]
例如 $x=-2$ 时仅 $t=-2$ 那项看到亮侧的 1，故 $g(-2)=-0.1$，出现欠冲；$x=1$ 时前四个系数都落在亮侧，故 $g(1)=-0.1+0.2+0.8+0.2=1.1$，出现过冲。真实理想低通具有更多旁瓣，所以可能反复交替形成多道波纹。

\begin{center}
\begin{tikzpicture}[x=.95cm,y=2.1cm,>=Latex,font=\small]
\draw[->] (-3.5,0)--(3.8,0) node[right]{$x$};
\draw[->] (-3.3,-.28)--(-3.3,1.4) node[above]{灰度};
\draw[gray,dashed] (-3,1)--(3.3,1);
\draw[thick,dashed] (-3,0)--(0,0)--(0,1)--(3,1);
\draw[very thick] plot coordinates{(-3,0)(-2,-.1)(-1,.1)(0,.9)(1,1.1)(2,1)(3,1)};
\foreach \x/\y in {-3/0,-2/-.1,-1/.1,0/.9,1/1.1,2/1,3/1}{\fill (\x,\y) circle(1.5pt);}
\draw[->] (-1.6,-.35)--(-2,-.12) node[pos=0,below]{低于暗侧};
\draw[->] (1.8,1.36)--(1,1.12) node[pos=0,above]{高于亮侧};
\node[align=left] at (4.4,.55){虚线：原阶跃\\实线：示意输出};
\end{tikzpicture}
\end{center}

作为对照，若使用非负、总和为 1 的核
\[
h_{\mathrm{pos}}=\frac1{16}[1,4,6,4,1],
\]
那么每个输出是邻域像素的加权平均，不会超出邻域最小值和最大值。对同一阶跃，过渡值依次为 $1/16,5/16,11/16,15/16,1$，单调地从 0 过渡到 1，没有过冲、欠冲。高斯空间核同样非负，这解释了常说的“高斯低通无振铃”。

这一结论指理想高斯核及恰当实现本身不因负旁瓣产生振铃；边界延拓、核截断、其他锐化步骤仍可能引入伪影。高斯也会让边缘过渡变宽，“没有晕圈”不等于“完全不模糊”。

\subsection{高通、带通、带阻和陷波：先判断想保留什么}

\subsubsection{高通不是自动得到一张好看的锐化图}
\textbf{高通滤波（high-pass filtering）}保留较快变化，压低较慢变化。本模块常用低通的补函数构造：
\[
H_{\mathrm{HP}}(u,v)=1-H_{\mathrm{LP}}(u,v).
\]
$H_{\mathrm{LP}}$ 是选定低通响应，$H_{\mathrm{HP}}$ 是对应高通响应。分别代入理想、巴特沃斯和高斯低通，就得到 IHPF、BHPF 和 GHPF，即三种对应的高通滤波器。例如某点低通为 0.8，高通就是 0.2；某点低通为 0.1，高通就是 0.9。低通在 DC 为 1，故这种高通在 DC 为 0，去掉整幅图的均值。

纯高通输出常包含正负值，主要表达边缘和细节，失去原来的整体亮度。若希望保留背景并增强细节，可以采用
\[
g=f+k\,g_{\mathrm{HP}},\qquad
H_{\mathrm{sharp}}=1+kH_{\mathrm{HP}},\qquad k>0.
\]
$g_{\mathrm{HP}}$ 是高通输出，$k$ 是加入细节的强度。以四像素例子为例，原图 $[12,10,8,10]$，纯高通为 $[2,0,-2,0]$，取 $k=0.5$，得到 $[13,10,7,10]$。平均仍为 10，局部起伏增大。

原稿说“高通放大噪声”，更准确的理解是：高频噪声会相对低频背景更突出；当使用 $0\leq H_{\mathrm{HP}}\leq1$ 的纯高通时，不是每个噪声频率的绝对振幅都增大。高提升响应 $1+kH_{\mathrm{HP}}$ 可大于 1，此时相应频率的信号和噪声都确实被放大。

\subsubsection{带阻与带通：按一整圈频率选择}
\textbf{带阻滤波（band-reject filtering）}压低某个频率范围，两侧通过。以径向理想带阻为例，
\[
H_{\mathrm{BR}}(D)=
\begin{cases}
0,&D_0-W/2\leq D\leq D_0+W/2,\\
1,&\text{其他}.
\end{cases}
\]
$D_0$ 是被抑制环带的中心半径，$W>0$ 是带宽，通常选 $D_0>W/2$ 以使内半径为正。例如 $D_0=20,W=6$，被抑制的是 $17\leq D\leq23$，不是“所有 $D>20$”。\textbf{带通滤波（band-pass filtering）}反过来只保留这个频带：$H_{\mathrm{BP}}=1-H_{\mathrm{BR}}$。

原稿给出的巴特沃斯带阻为
\[
H_{\mathrm{BR,B}}(D)=
\frac{1}{1+\left[\dfrac{DW}{D^2-D_0^2}\right]^{2n}}.
\]
$D,D_0,W$ 的含义同上，正整数阶数 $n$ 控制过渡陡峭程度。$D=D_0$ 处不能直接除以零，取极限规定 $H=0$。例如 $D_0=10,W=4,n=1,D=5$，括号内为 $20/(25-100)=-4/15$，平方后为 $16/225$，因此 $H=225/241\approx0.9336$。该点仍可通过约 $93.36\%$ 的振幅。

\begin{center}
\begin{tikzpicture}[x=1cm,y=1.5cm,font=\small,>=Latex]
\begin{scope}
\draw[->] (0,0)--(4.5,0) node[right]{$D$};
\draw[->] (0,0)--(0,1.3) node[above]{$H$};
\draw[thick] (0,1)--(1.5,1)--(1.5,0)--(2.8,0)--(2.8,1)--(4.2,1);
\node at (2,-.55){带阻：中间这段不通过};
\node[left] at (0,1){1};
\draw[<->] (1.5,.45)--(2.8,.45) node[midway,above]{$W$};
\end{scope}
\begin{scope}[xshift=7cm]
\draw[->] (0,0)--(4.5,0) node[right]{$D$};
\draw[->] (0,0)--(0,1.3) node[above]{$H$};
\draw[thick] (0,0)--(1.5,0)--(1.5,1)--(2.8,1)--(2.8,0)--(4.2,0);
\node at (2,-.55){带通：只让中间这段通过};
\node[left] at (0,1){1};
\end{scope}
\end{tikzpicture}
\end{center}

原稿还提到高斯型带阻：它也以平滑过渡替代突然的零一切换。理想、巴特沃斯、高斯这些名字描述的是过渡形状；带通、带阻描述的是保留哪个频段，两组分类不要混淆。

带阻可能消除某种周期干扰，但同时会删去同一半径上各个方向的真实纹理。若噪声只表现为两个局部亮点，抑制整圈通常过宽。

\subsubsection{陷波：只处理被确认的几个频率位置}
\textbf{陷波滤波（notch filtering）}是在特定频率位置附近开小孔。令中心化频率坐标为 $\boldsymbol q=(\Delta u,\Delta v)$，识别出的干扰位置为 $\boldsymbol q_0=(a,b)$。理想成对陷波可写成
\[
H_{\mathrm{notch}}(\boldsymbol q)=
\begin{cases}
0,&\|\boldsymbol q-\boldsymbol q_0\|\leq r
\ \text{或}\ \|\boldsymbol q+\boldsymbol q_0\|\leq r,\\
1,&\text{其他}.
\end{cases}
\]
$\|\cdot\|$ 表示欧氏距离，$r$ 是每个小孔的半径。例如干扰位于偏移 $(0,20)$ 与 $(0,-20)$，可以在两点附近各压低半径 2 的区域，而不把半径 20 的整圈都删掉。实际常用平滑陷波边缘，减少突截带来的振铃。

\begin{center}
\begin{tikzpicture}[font=\small]
\begin{scope}
\fill[gray!10] (-2,-2) rectangle(2,2);
\fill[gray!65,even odd rule] (0,0) circle(1.5) (0,0) circle(1.05);
\draw (-2,-2) rectangle(2,2);\draw[gray](-2,0)--(2,0);\draw[gray](0,-2)--(0,2);
\fill (0,0) circle(1.7pt);
\node[align=center] at (0,-2.55){环形带阻\\灰色环带被抑制};
\end{scope}
\begin{scope}[xshift=6.5cm]
\fill[gray!10] (-2,-2) rectangle(2,2);
\fill[gray!65] (-1.3,0) circle(.25);\fill[gray!65] (1.3,0) circle(.25);
\draw (-2,-2) rectangle(2,2);\draw[gray](-2,0)--(2,0);\draw[gray](0,-2)--(0,2);
\fill (0,0) circle(1.7pt);
\node[align=center] at (0,-2.55){成对陷波\\仅局部小孔被抑制};
\end{scope}
\end{tikzpicture}
\end{center}

对于实输入，要保持输出为实数，应让 $H$ 保持相应的共轭对称性。实数、偶对称的成对小孔是最方便的做法；只去掉一边通常破坏配对，使 IDFT 出现不能仅归因于舍入误差的虚部。DC 是自身的共轭伙伴，只令中心 $H=0$ 是一种特殊的直流陷波，作用是去均值，\textbf{不是通用的周期条纹去除方法}。

\subsubsection{噪声与高频并非同义词}
边缘和草地纹理是有用的高频信息；条纹干扰可能落在低频、中频或高频；白噪声通常在很宽的频率范围都有成分。图像上的缓慢阴影可能主要影响低频，也可能根本不是噪声。选择滤波器前应先问：干扰的频率在哪里？它与有用信号是否重叠？

若二者在同一频率位置重叠，简单频域乘法不能知道哪个是噪声、哪个是物体。抑制噪声也会损失信号。因此“外圈全部是噪声”“所有对称亮点都该删掉”“低通一定去噪且保边”都不成立。

\subsection{同态滤波：先解决“光照乘在表面上”这个问题}

\subsubsection{照明与反射为什么是相乘}
\textbf{同态滤波（homomorphic filtering）}从一个简化成像模型出发：
\[
f(x,y)=i(x,y)\,r(x,y).
\]
$f$ 是记录到的灰度，$i>0$ 是照明强度，$r$ 是表面反射比例；为使后面取对数有明确含义，$f$ 和 $i$ 均理解为选定参考尺度下的无量纲正数，$r$ 本身就是无量纲比例。在简化模型中常取 $0<r\leq1$。同一表面得到两倍的照明，记录亮度也大致变为两倍，这就是乘法的直观来源。例如 $i=100,r=0.3$ 时 $f=30$；照明升为 200，反射仍为 0.3，则 $f=60$。

常用假设是：照明在空间上变化较缓慢，反射在物体边界和纹理处变化较快。因此希望压低缓慢的大范围亮暗差，同时加强细节。\textbf{这只是近似：}阴影边界也可能很锐利，反射率也可能缓慢变化；两者的频率范围并不一定严格分开。

\subsubsection{为什么先取对数，而不是直接做高通}
直接对 $f=i\,r$ 变换，乘积对应频域中的卷积，照明与反射不会简单地成为两个可各自加权的频谱。对数有一个关键性质：$\ln(ab)=\ln a+\ln b$。所以令
\[
z(x,y)=\ln f(x,y)=\ln i(x,y)+\ln r(x,y).
\]
$z$ 是对数灰度。对它做 DFT，得到
\[
Z(u,v)=I_{\log}(u,v)+R_{\log}(u,v).
\]
$Z$ 是 $z$ 的频谱，$I_{\log}$ 是 $\ln i$ 的频谱，$R_{\log}$ 是 $\ln r$ 的频谱。现在两部分在数学上相加，频域滤波可写成
\[
S=HZ=H I_{\log}+H R_{\log}.
\]
$S$ 是滤波后的对数域频谱。同一个 $H$ 根据频率给予不同权重，在“照明偏低频、反射细节偏高频”的假设下起作用。\textbf{取对数并没有自动、精确地识别出照明和反射}；两者频谱若重叠，仍会相互影响。

数值上，若 $i=100,r=0.25$，则 $f=25$；$\ln100\approx4.60517$、$\ln0.25\approx-1.38629$，二者之和为 $3.21888\approx\ln25$。反射率的对数可能为负，这完全正常，不需要为此删掉负值。

\subsubsection{常用滤波器的每个参数是什么意思}
原稿采用
\[
H(D)=(\gamma_H-\gamma_L)
\left[1-\exp\!\left(-cD^2/D_0^2\right)\right]+\gamma_L.
\]
$D$ 是到频谱中心的距离；$D_0>0$ 是过渡尺度；$c>0$ 调节过渡速度；$\gamma_L$ 是低频端的增益；$\gamma_H$ 是高频端的增益。常见增强设置是 $0<\gamma_L<1<\gamma_H$。

先检查两端：当 $D=0$ 时，指数为 0，方括号为 $1-1=0$，故 $H(0)=\gamma_L$。当 $D$ 很大时，指数趋向很大的负数，指数函数趋近零，方括号趋向 1，故 $H(D)\to\gamma_H$。这正是“低频压低，高频提高”的来源。

例如取 $\gamma_L=0.5,\gamma_H=1.5,c=1$，则 $H(0)=0.5$；在 $D=D_0$ 处，
\[
H(D_0)=(1.5-0.5)(1-e^{-1})+0.5
\approx1.13212;
\]
远离中心时趋近 1.5。这里“低频减半”指\textbf{对数域分量的振幅减半}，不代表把原图每个暗区域的像素直接除以二。

\begin{center}
\begin{tikzpicture}[x=3.3cm,y=2cm,>=Latex,font=\small]
\draw[->] (0,0)--(2.8,0) node[right]{$D/D_0$};
\draw[->] (0,0)--(0,1.8) node[above]{$H$};
\draw[dotted] (0,.5)--(2.6,.5);\draw[dashed] (0,1.5)--(2.6,1.5);
\draw[thick,domain=0:2.6,samples=90] plot(\x,{1-exp(-\x^2)+.5});
\node[left] at (0,.5){$\gamma_L=0.5$};\node[left] at (0,1.5){$\gamma_H=1.5$};
\node[below] at (1,0){1};\draw[dotted] (1,0)--(1,1.1321);
\node[align=center,fill=white,inner sep=1pt] at (1.6,.38){压低缓慢变化\\提高快速变化};
\end{tikzpicture}
\end{center}

\subsubsection{完整流程与“动态范围压缩”的含义}
完整流程为：准备正值图像，取自然对数，按需要扩展边界与中心化，计算 DFT，乘 $H$，逆变换并正确移回，得到 $s(x,y)$，最后指数还原：
\[
g(x,y)=\exp\bigl(s(x,y)\bigr).
\]
$g$ 是输出图像，$s=\operatorname{IDFT}(S)$ 是滤波后的对数图。空间尺寸裁取和显示范围处理应与其他频域滤波一样明确说明。

如果某个正分量在对数域整体乘以 $\alpha$，指数还原后变成 $\exp(\alpha\ln i)=i^\alpha$。例如照明从 1 到 4，亮暗比为 4；取 $\alpha=0.5$ 后变成从 1 到 2，亮暗比降为 2，这就是乘法范围被压缩的直观含义。它不保证所有像素绝对变暗：对于 $0<i<1$，平方根反而比原值大。因此应说“压缩相对变化范围”，不能仅凭 $\gamma_L<1$ 就断言整幅图平均灰度必定下降。

\subsubsection{重要订正：两个增益相等，不等于普通线性滤波}
若 $\gamma_L=\gamma_H=\gamma$，前式第一项为零，$H(D)=\gamma$ 对所有频率都是常数。于是
\[
s=\gamma\ln f,\qquad g=e^{\gamma\ln f}=f^\gamma.
\]
当 $\gamma=1$ 时，输出等于输入；一般 $\gamma\neq1$ 时得到\textbf{幂律变换}，对原始灰度仍然非线性。例如 $\gamma=2$ 时，输入 2 和 3 分别输出 4 和 9，而输入和 5 输出 25，不等于 $4+9=13$，所以不满足线性的可加性。

原稿写成“$\gamma_L=\gamma_H$ 时退化为普通线性滤波”，需要改成：“在\textbf{对数域}退化为常数增益；在\textbf{原灰度域}通常是幂律非线性变换。”即使增益不同，整个“对数、线性滤波、指数”的组合一般也不是原图上的线性系统。

\subsubsection{零像素、噪声和模型限制}
$\ln0$ 没有有限值，不能直接对含零图像照搬公式。常见实现是先选尺度并使用 $z=\ln(\max(f,\varepsilon))$，其中 $\varepsilon>0$ 为小的正下限；这样明确改变了零附近的数据。另一种实现是 $\ln(1+f)$，但它不再严格满足 $\ln(ir)=\ln i+\ln r$ 的原模型，逆变换后要配套使用 $e^s-1$。

不能把 $\varepsilon$ 看成随便写的魔法数：它相对于灰度尺度多大，会影响暗部。接近零的值对微小噪声很敏感，高频提升也可能强化纹理噪声；过强参数可产生晕圈或不自然对比度。若图像严重饱和、照明与反射频谱高度重叠，简单同态模型不能保证恢复真实表面。

\subsection{原题 B4：频谱图判读，逐问建立推理链}
\label{original:B4}

\noindent\textbf{原题 B4（Enhancement3）}\quad
频谱图判读。(1) 中心亮点代表什么？(2) 远离中心的暗区代表什么？(3) 平坦的天空与密集的草地分别落在哪一部分？(4) 过中心出现一对对称亮斑，最可能是什么原因，用什么方法去除？

\paragraph{先确认题目的图是什么。}
本题通常指经过中心化、可能经过对数压缩显示的\textbf{幅度谱}。若图未中心化，DC 不在中央；若图显示的是相位，则“亮就是强”的读法也不成立。以下按中心化幅度谱作答。

\paragraph{(1) 中心亮点。}
中心是零频率，即直流分量，代表整体平均亮度。按照本模块 DFT 约定，中心系数的数值为全部像素之和，除以 $MN$ 得到平均值。中心\textbf{附近}的较亮区域还包括低频缓慢变化，不能把一大片中心区域都称作一个 DC 点。

\paragraph{(2) 远离中心的暗区。}
“远离中心”说明这里是较高的空间频率；“暗”说明这些频率的幅度较小。二者须分开解释。不能答成“原图这里很暗”，也不能只说“这里是高频，所以必然是噪声”。若幅度严格为零，则该离散频率分量缺失；显示为黑也可能仅是动态范围压缩后太弱而看不出来。

\paragraph{(3) 天空与草地。}
平坦天空在空间上变化慢，主要贡献 DC 和中心附近的低频；草地具有密集明暗变化，通常具有更多中高频成分。因此草地可能使外围某些区域增强。但草地也有平均亮度与大范围明暗，仍有低频；天空与地面的清晰分界也会贡献高频。不要说“一个物体只属于某一个频带”。

\paragraph{(4) 对称亮斑的原因与处理。}
如果原图同时可见不希望保留的等间距条纹，这对亮斑很可能来自周期干扰。实余弦以正、负频率配对，因此出现关于 DC 对称的亮点；两点连线垂直于条纹走向；离中心越远，条纹周期越短。

确认其为干扰后，优先在亮斑处使用\textbf{成对、局部、适度平滑的陷波}；若干扰能量形成完整环带，才考虑带阻。对称亮点也可能来自合法的栅格、布料或建筑纹理，所以仅凭频谱形状不能证明是噪声。频谱选点应结合原图和处理目标。

\paragraph{可整理成的短答案。}
中心为 DC，代表平均亮度；外围暗区为较弱的高频分量。平坦天空主要贡献低频，细密草地含较多中高频。对称亮点提示周期成分；若确认为条纹干扰，用成对陷波，必要时用带阻，并留意真实纹理损失。

\subsection{原题 D6：完整流程，重点是顺序与边界}
\label{original:D6}

\noindent\textbf{原题 D6}\quad
写出频域滤波的完整流程；并说明：(1) 为什么变换前要乘 $(-1)^{x+y}$？(2) 为什么补零到 $M\times N$？(3) 什么时候该用频域而不是 $3\times3$ 掩模？

\paragraph{先订正题目中尺寸的歧义。}
若题目的 $M\times N$ 指\textbf{最终选定的共同变换网格}，补到它是为让图像与核尺寸一致，但这个网格还必须足够大。若 $M\times N$ 指\textbf{原图尺寸}，则“核补到这个尺寸就一定得到线性卷积”不成立。下面统一用 $M\times N$ 表示原图、$P\times Q$ 表示变换网格。

\paragraph{完整流程：选择一条自洽的中心化路线。}
\begin{enumerate}[label=\arabic*.]
\item \textbf{明确目标与边界。}决定是平滑、提取细节还是去周期条纹；说明用零延拓、镜像还是周期边界。
\item \textbf{确定尺寸并准备数据。}若以 DFT 实现有限 $K\times L$ 核的完整线性卷积，选 $P\geq M+K-1,Q\geq N+L-1$，把图像和核填到共同尺寸，并处理核原点。
\item \textbf{正变换。}计算填充图像的 DFT，再用 \texttt{fftshift} 中心化；或者在偶数网格上先乘棋盘再做 DFT。二选一，不要把两条路线重复使用。
\item \textbf{建立 $H$ 并相乘。}让 $H$ 的中心、尺寸、频率单位与 $F$ 完全一致，逐点计算 $G=HF$。例如 $H=0.3$ 的频率位置保留原振幅的 $30\%$。
\item \textbf{逆变换并撤销移动。}使用频谱重排路线时先 \texttt{ifftshift} 再 IDFT；使用棋盘路线时先 IDFT 再乘一次棋盘。
\item \textbf{检查与截取。}对实输入和共轭对称滤波器，虚部应只剩微小舍入量；若虚部很大，先检查配对和索引。按核中心与输出需求裁取空间区域。
\item \textbf{单独处理显示范围。}若输出必须显示为 $0$ 到 $L_{\mathrm g}-1$ 的灰度，再选择缩放或截断；$L_{\mathrm g}$ 是可用灰度级数。保留负值的高通结果不宜未经说明就截断。
\end{enumerate}

这里“裁取到原图大小”是\textbf{空间尺寸}操作；“限制到灰度 $[0,L_{\mathrm g}-1]$”是\textbf{数值范围}操作。原稿把它们压在“裁剪”二字里，初学者需要明确区分。截断和归一化还会改变数值，不属于卷积定理本身。

\paragraph{(1) 乘棋盘的原因。}
在偶数网格中，乘 $(-1)^{x+y}$ 把 DC 平移到中心，方便用到中心的距离构造滤波器。它不是去噪，也不是增加高频。它只是可选的一种坐标安排，逆过程要乘回去；使用正确的 \texttt{fftshift}/\texttt{ifftshift} 也可以。

\paragraph{(2) 补零的原因与充分条件。}
DFT 乘法自然给出循环卷积。为求零延拓下的完整线性卷积，需要让两数组共同放入能容纳完整输出的网格。例：$256\times256$ 图像与 $3\times3$ 核至少用 $258\times258$。仅“尺寸一样”是可相乘的必要条件，\textbf{不等于}“尺寸足够大”这个避免折叠的条件。直接构造频域 $H$ 也仍需考虑周期边界。

\paragraph{(3) 什么时候频域更合适。}
$3\times3$ 固定核只需要每个像素约九次加权累加，空间域实现通常简单且开销小。大核卷积、需要精确选频的带阻与陷波、或者同态流程中的滤波，往往更适合考虑频域。

若图像为 $M\times N$、方核边长为 $k$，直接空间卷积粗略计算量为 $O(MNk^2)$；对 $P\times Q$ 网格做 FFT 的典型数量级为 $O(PQ\log(PQ))$。符号 $O$ 描述数据规模增大时的增长阶，不是实际耗时秒数。填充、缓存、硬件、可分离核和是否复用频谱都会影响速度，所以不存在不看实现便保证的统一分界。

\subsection{原题 D7：把同态滤波说清楚，避免四个跳步}
\label{original:D7}

\noindent\textbf{原题 D7}\quad
同态滤波。(1) 写出模型，说明照明与反射各对应低频还是高频；(2) 为什么必须先取对数？(3) 写出流程与 $H(u,v)$，并说明 $\gamma_L$、$\gamma_H$ 的作用。

\paragraph{(1) 模型与假设。}
写 $f(x,y)=i(x,y)r(x,y)$，再说明 $i$ 为照明、$r$ 为反射率。常假设照明变化慢，主要影响低频；反射的边界与纹理变化快，含较多高频。不能省略“常假设”，也不能声称反射率只有高频、照明绝无高频。

\paragraph{(2) 为什么取对数。}
原模型是相乘，而想利用的是频域中可加成分的线性加权。先取 $z=\ln f=\ln i+\ln r$，再变换成 $Z=I_{\log}+R_{\log}$，才能在这个加性表示中用同一 $H$ 对不同频率加不同增益。对数不保证两部分完全分离，只是让这种建模方式可用。含零像素要先选定正值预处理方法。

\paragraph{(3) 流程、公式和参数。}
流程写为
\[
f\ \xrightarrow{\ln}\ z\ \xrightarrow{\mathrm{DFT}}\ Z
\ \xrightarrow{\times H}\ S\ \xrightarrow{\mathrm{IDFT}}\ s
\ \xrightarrow{\exp}\ g.
\]
图中 $z$ 是对数图，$Z$ 是其频谱，$S$ 是滤波后的频谱，$s$ 是滤波后的对数图，$g$ 是输出；边界扩展、频谱移位、移回与裁取均应按前述规则配套执行。滤波器为
\[
H(u,v)=(\gamma_H-\gamma_L)
\left[1-e^{-cD(u,v)^2/D_0^2}\right]+\gamma_L.
\]
$D$ 为中心距离，$D_0,c$ 控制过渡；$\gamma_L<1$ 抑制对数域低频的相对变化，常用于减弱大范围照明差；$\gamma_H>1$ 提升对数域高频，常用于增强细节。由 $H(0)=\gamma_L$ 和远处 $H\to\gamma_H$ 验证，而不是只背两个不等式。

\paragraph{补充检验。}
令 $\gamma_L=\gamma_H=0.5$，得到 $g=f^{0.5}$。例如灰度 4 和 9 分别变成 2 和 3，说明这时是开平方的非线性变换，而不是普通线性低通。若两者同时为 1，才是这一设置下的恒等变换。

\subsection{原题 E6：逐步计算距离、增益与振铃}
\label{original:E6}

\noindent\textbf{原题 E6}\quad
频率域计算 + 振铃。$M=N=256$，频谱已中心化，$D_0=32$，某点偏移 $(\Delta u,\Delta v)=(24,32)$。(1) 求 $D$；(2) 求 ILPF、BLPF（$n=2$）、GLPF 在该点的 $H$；(3) 用 ILPF 平滑含边界的图像，为什么边界两侧出现一圈圈明暗环？哪种滤波器不振铃？

\paragraph{先读题：哪些数有用。}
$256\times256$ 告诉我们网格中心在 $(128,128)$。题目已经给出\textbf{相对中心的偏移} $(24,32)$，所以求距离时直接使用 24 和 32，不能再各减 128。若题目给的是绝对坐标 $(152,160)$，才先计算 $152-128=24$、$160-128=32$。

\paragraph{(1) 求半径。}
\[
\begin{aligned}
D&=\sqrt{(\Delta u)^2+(\Delta v)^2}\\
 &=\sqrt{24^2+32^2}=\sqrt{576+1024}
 =\sqrt{1600}=40.
\end{aligned}
\]
先平方、相加、开方，不能把距离算成 $24+32=56$。$D=40$ 与 $D_0=32$ 单位相同，均为频率格；其比值为 $D/D_0=40/32=5/4=1.25$。

\paragraph{(2a) 理想低通。}
它只需要比较内外：$40>32$，该点在通带圆盘外，因此
\[
\boxed{H_{\mathrm I}=0}.
\]
这意味着此频率系数被完全删除，不意味着原图某个像素被设为零。

\paragraph{(2b) 巴特沃斯低通，阶数为 2。}
先确定指数：$2n=2\times2=4$。然后代入，尽量到最后才舍入：
\[
\begin{aligned}
H_{\mathrm B}
&=\frac{1}{1+(D/D_0)^{2n}}
=\frac{1}{1+(40/32)^4}\\
&=\frac{1}{1+(5/4)^4}
=\frac{1}{1+625/256}
=\frac{256}{881}\\
&\approx\boxed{0.2905788876}.
\end{aligned}
\]
该频率保留约 $29.06\%$ 的\textbf{振幅}。若输入该点的幅度为 100，输出约为 $100\times0.2905788876=29.0579$；功率保留比例是这个数的平方，约 $0.08444$，不能把振幅百分比直接当成功率百分比。

\paragraph{(2c) 高斯低通。}
\[
\begin{aligned}
H_{\mathrm G}
&=\exp\!\left[-\frac{D^2}{2D_0^2}\right]
=\exp\!\left[-\frac{40^2}{2\times32^2}\right]\\
&=\exp\!\left[-\frac{1600}{2048}\right]
=\exp(-0.78125)\\
&\approx\boxed{0.4578333618}.
\end{aligned}
\]
同一频率保留约 $45.78\%$ 振幅。原稿的三位小数结果 $0.291$、$0.458$ 可以作为四舍五入答案；完整计算应保留上面的精确分数与指数，避免中间过早截断。

\begin{center}
\begin{tabular}{@{}llll@{}}
\toprule
滤波器&代入结果&约保留振幅&本点的解释\\
\midrule
ILPF&$0$&$0\%$&圆盘外完全去掉\\
BLPF，$n=2$&$256/881$&$29.06\%$&逐渐衰减\\
GLPF&$e^{-0.78125}$&$45.78\%$&平缓衰减\\
\bottomrule
\end{tabular}
\end{center}

本题不能据此断言“高斯在所有位置都比巴特沃斯保留更多”。比较只针对当前半径和参数；两条曲线在其他位置可能有不同的相对关系。

\paragraph{(3) 为什么产生明暗环，哪一种避免这种振铃。}
ILPF 在 $D_0$ 处突然从 1 变为 0；对应空间核具有振荡旁瓣，且包含负值。与阶跃边缘卷积时，一部分旁瓣跨过边界，把亮侧与暗侧按正负权重混合，产生过冲与欠冲。二维圆形低通可能在边界附近出现环状或平行于边界的明暗晕纹，这就是振铃。

在所比较的标准滤波器中，\textbf{高斯低通的理想空间核仍为非负高斯，没有负旁瓣，因此不会产生这种旁瓣振铃}；它仍会模糊边缘。巴特沃斯过渡较缓时通常振铃较轻，提高阶数趋近理想截止，振铃风险相应增大。实际实现还要检查边界处理。

\paragraph{本题的截止点纠错。}
$D=D_0$ 时，按本题公式，BLPF 始终为 $0.5$，对应振幅增益约 $-6.02$ dB；GLPF 为 $e^{-1/2}\approx0.60653$。原稿把 BLPF 的 $0.5$ 标成 $-3$ dB，是混用了振幅比、功率比或另一种巴特沃斯定义。

\subsection{做题前的检查表与自测}

\begin{enumerate}[label=\arabic*.]
\item \textbf{先问坐标。}看到的是像素图、幅度谱还是相位谱？是否中心化？横竖轴分别对应哪一方向？
\item \textbf{再问频率。}位置远近说明变化快慢，亮暗说明幅度强弱；不要把两者混为一谈。
\item \textbf{再问滤波目标。}平滑用低通；提取变化用高通；指定尺度用带通；环形干扰考虑带阻；局部成对周期分量考虑陷波。
\item \textbf{代入前看定义。}低通的 $D_0$ 不是跨所有公式通用的半功率点；BLPF 的幂是 $2n$，GLPF 的分母是 $2D_0^2$。
\item \textbf{逆变换前后看配套。}中心化要移回；图像和核要共同够大；核原点要对齐；虚部异常要检查，不能直接掩盖。
\item \textbf{最后看副作用。}平滑可能损失边缘；高提升可能放大噪声；硬截止可能振铃；对数模型需要正值且不保证照明与反射可完全分开。
\end{enumerate}

\paragraph{自测 1：}一幅图有间隔 8 像素的竖条纹，另一幅图间隔 4 像素。哪幅的亮点更远离中心？亮点连线大致朝哪里？\quad\textbf{答案：}4 像素周期的频率更高，亮点离中心更远；竖条纹沿水平方向发生变化，连线大致水平。比较需使用相同尺寸与采样单位。

\paragraph{自测 2：}某点在中心附近，但频谱显示很暗，能说“它是高频”吗？\quad\textbf{答案：}不能。它的位置说明频率低，暗说明该低频分量幅度小。

\paragraph{自测 3：}$100\times80$ 图像与 $5\times7$ 核求完整线性卷积，共同网格至少多大？\quad\textbf{答案：}$(100+5-1)\times(80+7-1)=104\times86$，图像和核都须放入该网格。若之后裁回原图大小，还要说明核中心与裁取位置。

\paragraph{自测 4：}同态滤波中 $\gamma_L=\gamma_H=2$，输出是什么？\quad\textbf{答案：}在正值、无额外偏移与显示缩放的模型下，$g=f^2$，不是原灰度域的线性滤波。若同时为 1，输出才保持不变。

\paragraph{这一模块的一句话主线。}\quad
先把图像理解为不同快慢、不同方向的亮度变化；DFT 给出这些变化的配方；$H$ 决定各成分保留多少；IDFT 把新配方还原成图像。频谱中心化负责坐标，补零与延拓负责边界，滤波器形状负责取舍，而相位、真实纹理与噪声重叠等限制决定了哪些信息不能凭一个简单滤波器恢复。
